\cotefile{7}
\scanpage{47}
non plus de théorème de dualité, ni de Künneth. \deleted{Il est d'autant \unsure{p|(|plus}}

(Il semble cependant que malgré cela, ils arrivent à utiliser leur cohomologie pour donner une expression à la Lefschetz-Weil de la fonction $\zeta$ d'un schéma projectif et lisse sur un corps fini, en montrant que \inserted{les} Frobenius opérant sur les $H^i(X)$ sont des opérateurs à traces .... ! Ce formalisme semble alors proche de celui de Dwork, que je ne connais pas .... ).

\S~3 \deleted{4)} \uline{Relations \inserted{\uline{heuristique de} W.M.} avec une idée de Gauss-Manin.}

1. Caractère assez surprenant à première vue de l'assertion d'invariance de W.M. Supposons p.ex. que l'on ait
% archive sheet number: 46

